% Chapter 4: Robust Task Policy via Observation Augmentation and Human-in-the-Loop Learning

\chapter{Robust Task Policy via Observation Augmentation and Human-in-the-Loop Learning}\label{ch:task_policy}

\section{Observation Space Design}\label{sec:obs_design}

% From POMDP analysis (Ch.3) to concrete 24D design:
%
% | Dim   | Content         | Source              | Biased? |
% |-------|-----------------|---------------------|---------|
% | 0-6   | Joint angles    | Encoder (q+b)       | Yes     |
% | 7-13  | Joint velocities| Encoder derivative  | No      |
% | 14    | Gripper         | Control command      | No      |
% | 15-17 | TCP position    | Biased FK(q+b)      | Yes     |
% | 18-20 | Block position  | External camera      | No      |
% | 21-23 | Real TCP        | External + 5mm noise | No      |
%
% Information-theoretic role of each component:
% - Biased signals: provide q+b information
% - Real TCP: provides (noisy) q information
% - Joint: both together enable implicit bias inference
% - Block pos: task goal, independent of bias

\section{Randomized Bias Training}\label{sec:random_bias}

% Motivation: treat b as environment parameter -> domain randomization
% b ~ U[0, b_max] per episode, constant within episode
%
% Comparison with fixed-bias training:
% - Fixed b=0.2: overfits to specific b -> overcompensates at b=0
% - Random b in [0, 0.25]: learns b-conditional policy pi(a|o, Delta_x)
%
% Connection to robust RL:
% - Not worst-case (robust MDP), but average-case over b distribution
% - Randomization forces policy to USE the Delta_x signal to infer b

\section{RLPD Training Framework}\label{sec:rlpd}

% SAC (Soft Actor-Critic) for continuous action space
%
% RLPD = SAC + Offline Demo Buffer + Online Buffer
% - Offline buffer: 50 human demonstration episodes
% - Online buffer: collected during training (100K capacity)
% - Warmup: 500 steps offline-only before mixed training
% - Replay mix: 50% offline + 50% online
%
% Human-in-the-loop:
% - Keyboard intervention during online collection
% - Successful interventions added to offline buffer
% - Failed interventions discarded
%
% Distributed architecture: Actor <-gRPC-> Learner

\section{Network Architecture}\label{sec:network}

% State encoder: Linear(24, 256) + LayerNorm + Tanh
% Image encoders: ResNet10 (frozen) for front/wrist cameras
% Actor: [256, 256] MLP -> 6D continuous + 1D discrete gripper
% Critic: 2x [256, 256] MLP (twin Q-learning)
% Temperature: learnable log_alpha, auto-tuned
%
% SAC losses:
% L_Q = E[(Q(s,a) - (r + gamma * min_i Q_target_i(s', a')))^2]
% L_pi = E[alpha * log pi(a|s) - Q(s, a)]
% L_alpha = -E[alpha * (log pi(a|s) + H_target)]
%
% Key hyperparameters:
% batch_size=256, utd_ratio=2, discount=0.97, lr=3e-4

\section{Reward Design}\label{sec:task_reward}

% Sparse reward for pick-cube:
% r = 1.0 if success else 0.0
% success = (distance < 5cm) AND (height_gain > 20cm)
%
% Why sparse works here:
% - RLPD provides demo-guided exploration
% - Human interventions bootstrap early training
% - Dense reward risks reward hacking under bias
